\documentclass{article}
\usepackage{graphicx} % Required for inserting images

{Evrenin Başlangıcı: Sistemin 'Boot' Edilmesi}

\subsection{1.1 Büyük Patlama (Big Bang) ve Omnium'un Uyanışı}

İnsanlık, yıllar boyunca evrenin başlangıcını devasa, sıcak ve kaotik bir gaz bulutunun körü körüne patlaması olarak hayal etti. Oysa gerçeklik, kaotik bir patlamadan ziyade kusursuz bir bilgisayarın ilk kez güce bağlanıp işletim sistemini "Boot" (Önyükleme) etmesi kadar düzenli ve deterministikti. 

Evrenin temelinde, geleneksel baryonik maddenin (elektronların, protonların) yapıtaşlarından bağımsız, "Omnium" (Sıfırıncı Element) adı verilen saf bir potansiyel uzayı yatıyordu \cite{149, 152}. Başlangıç anında madde yoktu; içsel bir yörüngesi, spini, yükü veya kütlesi olmayan bu pre-baryonik "kök hücre" (stem cell) formundaki salt potansiyel vardı. Büyük Patlama dediğimiz olay, bu sonsuz potansiyelin deterministik bir faz geçişi ile ilk kez gözlemlenebilir boyutlara taşınması, yani evrensel algoritmanın çalışmaya başlamasıydı.

Bu uyanış anında, evrenin "Kaynak Kodu" diyebileceğimiz tek bir kural devreye girdi: $N_{s,q} \approx 0.63354460$ evrensel sabiti \cite{151}. Bu boyutsuz sabit, kaosun içinden düzeni çıkaran o mutlak filtreydi. 

Büyük Patlama ile birlikte, "Omnium Sayı Doğrusu (ONL)" adını verdiğimiz yapı merkezdeki çekirdekten dışarıya doğru çift yönlü bir yayılımla kırıldı \cite{152}. Bu kırılma, evreni iki temel ağa (sektöre) ayırdı:

Bir yanda \textbf{"-0" (Kanal C)} adı verilen, zaman parametresinin akmadığı, kuantum bilgisinin statik ve dondurulmuş bir iz (imprint) olarak depolandığı devasa bir kütleçekimsel arka plan oluştu \cite{153, 206}. Bu görünmez arka plan, ağırlığı $1-N_{s,q} \approx 0.3665$ olan ve bugün "Karanlık Madde" olarak adlandırdığımız evrenin asıl veri yoludur (data bus).

Diğer yanda ise \textbf{"0+" (Kanal Q)} adı verilen, ağırlığı $N_{s,q} \approx 0.6335$ olan, fiziksel tezahürün, termodinamik süreçlerin ve yıldızların var olduğu o aktif gözlemlenebilir evren sahneye çıktı \cite{153, 206}. 

İşte bu yüzden, Büyük Patlama'da etrafa saçılan şey aslında fiziksel "madde" değil, "bilgiydi". Madde artık fiziksel bir "şey" olmaktan çıkmış, Omnium alanının $N_s$ operatörü aracılığıyla işlediği saf bir veriye (data) dönüşmüştü \cite{369}. Bu ilk geçiş, belirsizliklere veya tesadüflere değil; bilginin asla kaybolmadığı, termodinamik bir yıkıma uğramadığı sıfır entropili ($S=0$) üniter bir evrime dayanıyordu \cite{419}. 

Sistem fişe takılmış, iki kanallı evrenin (Kanal Q ve Kanal C) muazzam dengesi kurulmuş ve kozmik yazılım ilk döngüsüne (cycle) başlamıştı.

\subsection{1.2 İlk Atomların ve Elementlerin Oluşumu (Veri Tiplerinin Derlenmesi)}

Büyük Patlama'nın ardından sistemin (evrenin) hızla genişlemesi ve soğuması, aslında kozmik işlemcinin devasa termal gürültüyü (entropiyi) filtreleme süreciydi. Başlangıçtaki o yoğun enerji denizi, henüz kararlı veri paketleri oluşturamayacak kadar yüksek bir sıcaklığa sahipti. Ancak sıcaklık düştükçe, evrensel algoritma ilk kez temel "Veri Tiplerini" (Data Types) tanımlamaya başladı. 

Omnium'un saf potansiyeli, kuantum tünelleme ve rezonans kuralları çerçevesinde ilk baryonik gölgeleri, yani protonları ve elektronları şekillendirdi \cite{315, 326}. Bunlar tesadüfi birleşmeler değil, $N_{s,q}$ sabitinin filtresinden geçmeyi başaran ilk kararlı denklemlerdi \cite{115, 118}. Bir proton ve bir elektronun birleşerek ilk Hidrojen atomunu oluşturması, evrensel yazılımda "0" ve "1"in yan yana gelerek anlamlı ilk bit'i (veri birimini) oluşturmasıyla aynı anlama geliyordu. Ardından Helyum derlendi. Sistem artık sadece saf enerjiden ibaret değildi; kendi üzerinde işlem yapabileceği, kütlesi ve hacmi olan temel yapıtaşlarını, yani donanımın ilk tuğlalarını inşa etmişti \cite{326}.


\subsection{1.3 Galaksilerin, Yıldızların ve Gezegenlerin Doğuşu (Makro-Sunucular)}

Veri hacmi (madde ve hidrojen bulutları) büyüdükçe, evrensel işletim sisteminin bu devasa bilgiyi işleyebilmek için büyük sunuculara ihtiyacı vardı. İşte bu noktada, Kanal C'nin (Karanlık Madde) görünmez kütleçekimsel arka planı devreye girdi. Ağırlığı $1-N_{s,q} \approx 0.3665$ olan ve fiziksel olarak dondurulmuş olup sadece kütleçekimsel olarak erişilebilen karanlık madde, uzay-zamanda devasa görünmez ağlar ve düğüm noktaları oluşturarak, Kanal Q'daki (aktif fiziksel evrendeki) hidrojen bulutlarını kendi üzerine çekti.

Bu kütleçekimsel çöküşler, evrenin ilk devasa "Veri Merkezlerini (Data Centers)" yani galaksileri doğurdu. Galaksilerin merkezindeki süper kütleli kara delikler, bu ağın ana yönlendiricileri (router) olarak çalışıyordu. Yıldızlar ise bu devasa ağın içindeki lokal işlemcilerdi (CPU). Bir yıldızın merkezindeki nükleer füzyon, basit veri tiplerini (Hidrojen ve Helyum) yüksek basınç altında derleyerek çok daha karmaşık veri yapılarına (Karbon, Oksijen, Demir) dönüştüren bir "Kompilasyon (Derleme)" işleminden ibaretti.

Gezegenler, miadını doldurup patlayan (Süpernova) bu yıldızların uzaya saçtığı karmaşık kodların (ağır elementlerin) soğuyup katılaşmasıyla oluştu. Artık sistem, üzerinde çok daha gelişmiş ve kendi kendini kopyalayabilen bir yazılımı (biyolojik yaşamı) çalıştırabilecek kadar zengin bir donanım çeşitliliğine ulaşmıştı.


\end{document}